% Generated draft: requires mathematical and source review.
\documentclass[11pt]{article}
\usepackage[margin=25mm]{geometry}
\usepackage{fontspec}
\setmainfont{Latin Modern Roman}
\usepackage{amsmath,amssymb,mathtools,unicode-math}
\setmathfont{Latin Modern Math}
\usepackage{enumitem,needspace,xurl,hyperref}
\setlength{\parindent}{0pt}
\setlength{\parskip}{5pt}
\setlength{\emergencystretch}{3em}
\newcommand{\N}{\mathbb{N}}
\newcommand{\Z}{\mathbb{Z}}
\newcommand{\Q}{\mathbb{Q}}
\newcommand{\R}{\mathbb{R}}
\newcommand{\C}{\mathbb{C}}
\newcommand{\F}{\mathbb{F}}
\newcommand{\A}{\mathbb{A}}
\newcommand{\PP}{\mathbb{P}}
\newcommand{\E}{\mathbb{E}}
\newcommand{\eps}{\varepsilon}
\newcommand{\dd}{\,\mathrm{d}}
\begin{document}
\section*{UnsolvedMath ID: 800991}
\textbf{Source problem code:} catalogue entry
\subsection{Definitions and mathematical statement}
Let $\emptyset'$ denote the halting set, and let $\mathrm{ML}(\emptyset')$ denote the class of reals that are Martin--L\"of random relative to $\emptyset'$, equivalently the $2$-random reals. In the usual binary-sequence presentation, randomness relative to $\emptyset'$ means avoiding $\bigcap_{n\in\N}U_n$ for every sequence $(U_n)_{n\in\N}$ of uniformly $\emptyset'$-effectively open subsets of $2^{\N}$ with $\mu(U_n)\leq 2^{-n}$, where $\mu$ is the fair-coin product measure. This explains the standard randomness terminology; the supplied excerpt does not specify its representation of reals.

For a specified comparison relation $\preccurlyeq$, a subset $B$ is a chain if
\[
\forall x,y\in B\quad (x\preccurlyeq y\ \lor\ y\preccurlyeq x).
\]
The supplied question is whether there exists a chain $B$ such that
\[
B\subseteq\mathrm{ML}(\emptyset')
\quad\text{and}\quad
\forall\alpha\in\mathrm{ML}(\emptyset')\quad [B]\neq[\alpha].
\]
Here $\preccurlyeq$ is only a placeholder for the intended comparison relation. The supplied record does not define that relation, the construction $[B]$, or the construction $[\alpha]$. In particular, it does not specify the objects being compared by the displayed inequality. Consequently, this display preserves the source question but is not a fully determined mathematical assertion until those definitions are recovered. No nonemptiness, cardinality, or effectiveness requirement on $B$ is supplied. The introductory description of this class as a special case refers to surrounding material absent from the record.
\subsection{Short English statement}
The source asks whether some chain of $2$-random reals has an associated object $[B]$ different from $[\alpha]$ for every individual $2$-random real $\alpha$. The ordering and bracket constructions are missing, so the intended statement requires clarification from the source context.
\subsection{Sources}
\begin{itemize}
\item Supplied local UnsolvedMath registry record: permanent ID $800991$, exact source problem code \texttt{catalogue entry}, title ``Open Problems in Mathematical Logic --- Problem 3.'' Upstream revision: \texttt{d371ee917ac9647fbb06cd6fc529deec6f3dda82}.
\item The record attributes the problem to Barmpalias, Gao et al., with proposed year $2026$, and identifies Problem 3 of \emph{Open Problems in Mathematical Logic} (2026): \url{https://arxiv.org/abs/2608.26628}. The supplied extraction method is \texttt{arxiv\_latex}; the supplied access date is 2026-10-16. These are supplied provenance claims, not independent source verification.
\item The local record labels the problem \emph{open} and records its update at \texttt{2026-10-16T14:00:00Z}. This status is supplied registry curation, not an independently verified mathematical or literature assessment.
\end{itemize}
\end{document}
